\section{Spesifisitas dan Cascade (Review Lanjut)}

Cascade dan specificity adalah mekanisme inti CSS yang menentukan aturan mana yang ``menang'' ketika beberapa selector menargetkan elemen yang sama. Meskipun sudah diperkenalkan di Bab 6, pemahaman mendalam diperlukan saat proyek CSS menjadi besar dan konflik aturan sering terjadi. Kegagalan memahami cascade sering menyebabkan penggunaan \code{!important} yang berlebihan \cite{mdnCss}.

\begin{table}[htbp]
\centering
\begin{tabular}{|P{3cm}|P{3.5cm}|P{2cm}|P{3.5cm}|}
\hline
\textbf{Level} & \textbf{Sumber} & \textbf{Prioritas} & \textbf{Contoh} \\
\hline
Inline & Atribut \code{style} & Tertinggi & \code{style="color:red"} \\
\hline
ID selectors & \code{\#id} & Tinggi & \code{\#nav \{ ... \}} \\
\hline
Class/pseudo & \code{.class}, \code{:hover} & Sedang & \code{.btn:hover} \\
\hline
Element/pseudo-el & \code{p}, \code{::before} & Rendah & \code{p \{ ... \}} \\
\hline
\code{!important} & Apapun & Override semua & Hindari jika mungkin \\
\hline
\end{tabular}
\caption{Hierarki specificity CSS (dari tinggi ke rendah)}
\end{table}

Praktik baik untuk mengelola specificity: (1) gunakan class sebagai selector utama (specificity konsisten); (2) hindari nesting selector terlalu dalam (\code{nav ul li a}---specificity tinggi, sulit di-override); (3) hindari ID selector untuk styling (specificity terlalu tinggi); (4) gunakan \code{!important} hanya sebagai pilihan terakhir; (5) pertimbangkan menggunakan CSS Layers (\code{@layer}) untuk mengontrol cascade di proyek besar \cite{webdevCss}.

\begin{konsep}
\textbf{CSS Layers (\code{@layer}).} Fitur baru CSS Cascade Layers memungkinkan pengelompokan aturan CSS ke dalam layer yang memiliki urutan prioritas eksplisit: \code{@layer base, components, utilities;}. Aturan di layer terakhir menang atas layer sebelumnya, terlepas dari specificity selector. Ini mengatasi masalah specificity war di proyek besar dan mempermudah integrasi library CSS pihak ketiga \cite{w3cCss}.
\end{konsep}

